% V2第19章思考题19.1--19.11，原扫描PDF185--186；图19.21--19.24需核验。
\section*{第19章\quad 磁场中的磁介质：思考题解析}
\paragraph{思考题 19.1}
(1) 在磁静态下 \(\nabla\times\boldsymbol H=\boldsymbol J_{\rm free}\) 的\emph{环流}只由自由电流决定，但局部H还取决于磁化强度和边界，因此“大小只由电流而与材料无关”不严格。(2) 在均匀、各向同性且 \(\mu>0\) 的线性顺磁/抗磁介质有 \(\boldsymbol B=\mu\boldsymbol H\) 同向；各向异性材料可不完全共线。(3) \(\oint\boldsymbol B\cdot d\boldsymbol S=0\) 恒成立。(4) \(\nabla\cdot\boldsymbol H=-\nabla\cdot\boldsymbol M\)，包围磁化端部时 \(\oint\boldsymbol H\cdot d\boldsymbol S\) 可以非零。因此不能把磁通B和磁化辅助场H的高斯定理混同。

\paragraph{思考题 19.2}
非均匀磁场中的小体积线性磁质受力数量级 \(\boldsymbol F\propto \chi\nabla B^2\)（弱磁化近似）。若线圈内部下方磁场更强，\(\nabla B^2\) 向下：顺磁质 \(\chi>0\) 被\emph{吸入强场向下}，抗磁质 \(\chi<0\) 被\emph{推出强场向上}。因此在图19.21该场梯度条件下，\(\boxed{\text{向上受力为抗磁质}}\)。若图中实际更强场位置相反，判断也需反向。

\paragraph{思考题 19.3}
无磁单极：任意闭面 \(\boxed{\oint\boldsymbol B\cdot d\boldsymbol S=0}\)，即使只围住永磁体的“北极”形状。因为 \(\boldsymbol H=\boldsymbol B/\mu_0-\boldsymbol M\)，\(\boxed{\oint\boldsymbol H\cdot d\boldsymbol S=-\oint\boldsymbol M\cdot d\boldsymbol S}\)；该面穿过磁体并切断磁化方向时H净通量一般不为零。这说明所谓有效磁荷只是在H场中的表面/体磁化源，不是真实B磁单极。

\paragraph{思考题 19.4}
铁磁材料具有磁畴和磁滞；敲击、加热或振动使磁畴壁跨过缺陷，残余磁化可能减弱。若同时施加地磁场并沿南北取向，外磁场可偏置磁畴的重排，使棒形成弱磁化；具体效果受矫顽力与合金性质影响。不能断言一次敲击必使所有磁体退磁或磁化。

\paragraph{思考题 19.5}
未磁化铁块的磁畴宏观磁矩之和近零，但外加永久磁铁的不均匀磁场会使铁块\emph{感应磁化}，在靠近磁铁的端部产生相反有效极性，且力近似指向更强场区，\(F\propto\chi\nabla B^2\)（线性弱场近似）。因铁磁响应还具有非线性与饱和，这只是定性模型。

\paragraph{思考题 19.6}
衔铁跨接马蹄形磁铁两极，减少空气隙磁阻及开放磁极的退磁场，从而有利于保存剩磁。多根条形磁铁N、S交错并列也可让磁通更充分闭合并减少整体退磁作用。此为传统永磁体保存措施；现代高矫顽力稀土磁体的最佳储存仍须按材料与安全要求决定。

\paragraph{思考题 19.7}
顺磁质磁矩定向的热无序竞争常给居里定律 \(\chi=C/T\)；铁磁质还存在畴结构和集体相变，其磁导率在居里温度附近尤其敏感且非线性。许多抗磁质的磁矩由外场改变电子轨道的感生电流决定，不依赖已有取向磁矩的热排序，故在普通温区 \(\chi_{\rm dia}\) 温度依赖较弱。不能把“顺磁”与“铁磁”的温度规律完全等同。

\paragraph{思考题 19.8}
由 \(\nabla\cdot\boldsymbol B=0\)，在没有漏磁的磁路任一管状支路横截面的磁通 \(\Phi=\int\boldsymbol B\cdot d\boldsymbol S\) 不随沿支路位置改变；在磁路分叉处，同一闭合面需满足 \(\boxed{\Phi_{\rm in}=\sum_i\Phi_{{\rm out},i}}\)。这类似电荷守恒的 Kirchhoff 电流定律。磁势降满足 \(\oint\boldsymbol H\cdot d\boldsymbol l=NI\) 且可用磁阻 \(\mathcal R=l/(\mu A)\)，但磁路是\emph{近似模型}，需要控制漏磁与边缘场。

\paragraph{思考题 19.9}
先在与冷源接触时施加强磁场，使顺磁离子的磁矩趋向有序并向冷源放出磁熵降低时的热；随后\emph{热隔离}，缓慢撤去外磁场，磁矩重新趋向无序，增加的磁熵由晶格熵减少补偿，导致温度进一步下降。理想绝热过程近似 \(S(B,T)=\mathrm{const}\)。要求低温下磁自由度仍有可调熵，残余内部相互作用和低温磁有序会限制冷却下限。

\paragraph{思考题 19.10}
永磁磁针在水平地磁场作用下趋向磁场水平分量方向；地理北极附近与之吸引的实际地磁极磁性相当于\emph{磁南极}，因此磁针“指北端”的磁学命名与地球磁极的地理位置不能混淆。磁偏角是磁北方向与真北方向之差，随时间、地理位置变化；在文献中“正南”“正北”还取决于历史标记法与观测基准。

\paragraph{思考题 19.11}
由界面边界条件 \(B_{1n}=B_{2n}\)（无磁单极），若铁芯和气隙的法向横截面积相同、漏磁/边缘场可忽略，则其\emph{平均法向磁通密度}近似相等。但铁芯端与空气隙边缘磁场不均匀时局部B不必完全相等。磁路近似 \(\boxed{\Phi\simeq NI/(\mathcal R_{\rm iron}+\mathcal R_{\rm gap})}\)，气隙越宽磁阻 \(\mathcal R_{\rm gap}\simeq l_g/(\mu_0A)\) 越大，所以同NI下宽气隙一般产生\emph{更小}的B；两种宽度不能视为相同。
